\section*{Chapitre \ref{ch:m2:emerging-market-fx-and-ndfs} --- Change des marchés émergents et NDF}

\begin{solution}{exo:m2:emerging-market-fx-and-ndfs:1}
Le contraire de l'acheteur~: il verse USD~319\,186 à 1\,425.50 et reçoit
USD~222\,222 à 1\,350.00.
\end{solution}

\begin{solution}{exo:m2:emerging-market-fx-and-ndfs:2}
Les roupies ne peuvent pas être livrées sur le \omterm{def:m2:emerging-market-fx-and-ndfs:onshore}{marché offshore}, et la couverture sur le \omterm{def:m2:emerging-market-fx-and-ndfs:onshore}{marché onshore} est restreinte par les
contrôles sur qui peut négocier et pourquoi. Le NDF se règle en dollars hors d'Inde
sur la base du fixing publié, si bien que le fonds peut couvrir sans déplacer de roupies.
\end{solution}

\begin{solution}{exo:m2:emerging-market-fx-and-ndfs:3}
6.9580 à 7.2420, 2~\% de part et d'autre de 7.1000~; le dollar de Hong Kong, 7.75 à 7.85
pour un dollar américain.
\end{solution}

\begin{solution}{exo:m2:emerging-market-fx-and-ndfs:4}
Deux jours ouvrés avant le lundi 5 octobre, en sautant le week-end et le vendredi
férié~: mercredi 30 septembre 2026.
\end{solution}

\begin{solution}{exo:m2:emerging-market-fx-and-ndfs:5}
Les contrats à terme offshore impliquent un taux du renminbi plus élevé que les contrats onshore. Un investisseur
qui couvre des obligations en renminbi en vendant du renminbi à terme perçoit le taux du dollar moins
du taux implicite du renminbi sur la couverture~: par rapport au comptant, la prime à terme
du renminbi est plus faible sur le \omterm{def:m2:emerging-market-fx-and-ndfs:onshore}{marché offshore} (7.12 à 7.07) que sur le \omterm{def:m2:emerging-market-fx-and-ndfs:onshore}{marché onshore} (7.10 à 7.04), si bien que
la couverture offshore rapporte donc moins de dollars. Elle coûte environ une base supplémentaire, et les \omterm{def:m2:emerging-market-fx-and-ndfs:controls}{contrôles des capitaux}
empêchent l'investisseur d'utiliser le cours onshore.
\end{solution}

\begin{solution}{exo:m2:emerging-market-fx-and-ndfs:6}
Pour empêcher sa monnaie de s'apprécier, la banque centrale vend sa propre monnaie,
qu'elle peut créer sans limite~; elle accumule des réserves de change, un risque pour son bilan
mais non une contrainte. Pour empêcher sa monnaie de se déprécier, elle doit vendre des devises étrangères,
dont les réserves s'épuisent.
\end{solution}

\begin{solution}{exo:m2:emerging-market-fx-and-ndfs:7}
Le taux de référence le plus bas tant que le plancher a tenu était de 1.2008, le 1\up{er} juin 2012. Il
valait 1.2010 le 14 janvier 2015 et 1.0280 le 15 janvier~; le plus bas de janvier était
0.9816, le 23 janvier.
\end{solution}

\begin{solution}{exo:m2:emerging-market-fx-and-ndfs:8}
La faible volatilité mesure la confiance du marché dans le plancher, et non l'ampleur
du mouvement en cas de rupture~; le risque est un saut, que l'ordre stop ne limite pas. Un ordre stop
devient un ordre au marché et s'exécute au cours suivant, qui le 15 janvier 2015 se situait
plus de 14~\% inférieur au taux de référence de l'après-midi. Le risque de la position était
le produit de la suppression du plancher par le saut, et non 2~\%.
\end{solution}

\begin{solution}{pb:m2:emerging-market-fx-and-ndfs:1}
\textbf{1.} $1\,000\,000 \times 1.2010/20 = $ CHF~60\,050.
\textbf{2.} $1\,000\,000 \times (1.2010 - 1.19) = $ CHF~11\,000, si l'ordre avait été exécuté à
son niveau.
\textbf{3.} $1 - 1.0280/1.2010 = 14.4\%$.
\textbf{4.} CHF~173\,000, soit 2.88 fois la marge.
\textbf{5.} $-$CHF~112\,950~: le client doit ce montant au courtier.
\textbf{6.} Un ordre stop de vente déclenché lorsque le cours atteint ou passe sous son niveau~; il devient alors
un ordre au marché, qui garantit l'exécution et non le prix.
\textbf{7.} Tous ceux qui s'étaient appuyés sur le plancher voulaient vendre des euros en même temps, et
les acheteurs se sont retirés le temps de juger où le taux allait se stabiliser.
\textbf{8.} Des années de cotations juste au-dessus de 1.20 faisaient paraître la baisse plafonnée~; le
levier avait été dimensionné pour la volatilité observée, et non pour la fin de la politique.
\textbf{9.} Faible, reflet de la crédibilité du plancher~; elle ne mesurait pas l'ampleur du
saut si le plancher venait à céder, car ce risque n'est pas continu.
\textbf{10.} Une option de vente de prix d'exercice proche de 1.20 coûtait peu avant et rapportait environ 17 centimes
par euro cet après-midi-là~: elle aurait couvert la perte là où l'ordre stop a échoué.
\textbf{11.} Le courtier, qui doit à ses propres contreparties les pertes sur les positions
perdantes du client, que celui-ci paie ou non.
\textbf{12.} Un prêt de deux ans de USD~300 millions consenti par Leucadia, conclu le 16
janvier 2015, dont le produit a remplacé le capital de ses entités réglementées.
\textbf{13.} Un taux initial de 10~\% par an, majoré de 1.5 point par trimestre, pour atteindre au plus
17~\%, avec des covenants et des remboursements anticipés obligatoires.
\textbf{14.} Par le saut que pourrait provoquer la fin de la politique, et non par la volatilité récente,
et avec une marge qui le couvre.
\textbf{15.} La protection contre le solde débiteur négatif, que le régulateur impose aux
courtiers de détail pour les produits à effet de levier~: les pertes du client sont plafonnées au dépôt.
\textbf{16.} Une sortie progressive serait anticipée et attaquée~; une surprise évite
d'acheter des devises étrangères sans limite dans les jours précédents, au prix d'un saut.
\textbf{17.} Tout engagement unilatéral~: d'autres planchers et plafonds de change, des \omterm{def:m2:emerging-market-fx-and-ndfs:peg}{ancrages de change}
défendus avec des réserves limitées, et des \omterm{def:m2:emerging-market-fx-and-ndfs:peg}{bandes de fluctuation} gérées près de leurs limites.
\textbf{18.} Par la perte si l'ancrage cède, dimensionnée à partir du saut et du capital
que l'on peut perdre, et non à partir de la volatilité du taux ancré.
\textbf{19.} Résultat à retenir~: \emph{le plancher qui a cédé} a coûté à la position longue à 20 pour 1
2.88 fois sa marge, CHF~173\,000 sur EUR~1 million au taux de référence de la BCE du
15 janvier 2015, ce qui laisse au client une dette de CHF~112\,950.
\textbf{20.} La volatilité reste faible jusqu'à la fin de la politique, puis elle bondit.
\end{solution}

\begin{solution}{iq:m2:emerging-market-fx-and-ndfs:1}
Un contrat à terme sur une paire de devises réglé uniquement en dollars, sur la différence entre
le taux du contrat et un fixing publié deux jours ouvrés avant le règlement-livraison~:
l'acheteur de dollars reçoit $N(S-K)/S$. Il existe pour les devises qui ne peuvent pas être
livrées sur le \omterm{def:m2:emerging-market-fx-and-ndfs:onshore}{marché offshore}.
\iqlookfor{règlement en espèces sur un fixing nommé, et raison de son existence.}
\end{solution}

\begin{solution}{iq:m2:emerging-market-fx-and-ndfs:2}
Les \omterm{def:m2:emerging-market-fx-and-ndfs:controls}{contrôles des capitaux} séparent les marchés~: les capitaux ne peuvent pas circuler librement entre
la Chine continentale et Hong Kong, si bien que l'offre et la demande sur les \omterm{def:m2:emerging-market-fx-and-ndfs:onshore}{marchés onshore} et offshore peuvent
différer, et l'arbitrage par les canaux autorisés est limité dans son ampleur. La
différence s'élargit lorsque les flux sont unilatéraux et que les canaux sont saturés.
\iqlookfor{les contrôles et les limites à l'arbitrage.}
\end{solution}

\begin{solution}{iq:m2:emerging-market-fx-and-ndfs:3}
Sous forme de mélange~: une diffusion autour de l'ancrage avec une faible volatilité, plus un saut de
taille incertaine à un instant aléatoire dont l'intensité reflète la
crédibilité de la politique (diffusion à sauts). Calibrer l'intensité et la taille des sauts sur les
risk reversals et les options très hors de la monnaie, qui valorisent la rupture~; la volatilité à la monnaie
en dit peu.
\iqlookfor{une composante de saut, et d'où viennent ses paramètres.}
\end{solution}

\begin{solution}{iq:m2:emerging-market-fx-and-ndfs:4}
Les positions dimensionnées pour un marché calme et plafonné ont subi un saut de 14~\%~; les ordres stop
se sont exécutés très en dessous de leurs niveaux~; les clients à effet de levier ont perdu plus que leur dépôt, et les courtiers qui
ne pouvaient pas recouvrer les soldes débiteurs ont dû lever du capital ou ont fait faillite. Les modèles de risque
avaient retenu la volatilité récente plutôt que l'événement.
\iqlookfor{risque de saut, ordres stop non exécutés, effet de levier et perte de contrepartie.}
\end{solution}

\begin{solution}{iq:m2:emerging-market-fx-and-ndfs:5}
Vendre des reals à terme contre des dollars au moyen de NDF sur le fixing PTAX, reconduits ou
appariés à l'horizon des obligations~; il faut vérifier la tarification onshore face à offshore et la
base. Au moment du fixing~: retards ou perturbations de publication (les mécanismes de repli standard
s'appliquent), un fixing qui diffère du taux auquel les obligations sont vendues, et
la base qui évolue contre la couverture.
\iqlookfor{NDF sur le fixing nommé, base et risque de fixing.}
\end{solution}

\begin{solution}{iq:m2:emerging-market-fx-and-ndfs:6}
Marges calculées sur des scénarios de stress autant que sur la volatilité~: sauts sur les ancrages et après
des événements connus, appliquées paire par paire~; plafonds d'effet de levier par client et par paire~;
limites de concentration sur les positions très répandues~; couvrir le portefeuille agrégé des clients
auprès de fournisseurs de liquidité et modéliser leurs propres risques de gap~; capital dimensionné pour
le pire scénario, soit les soldes débiteurs irrécouvrables~; la protection contre le solde négatif
intégrée dans les prix~; surveillance en temps réel avec réduction automatique de l'exposition avant les
événements programmés.
\iqlookfor{marge par scénario, plafonds d'effet de levier et fonds propres pour le risque de gap.}
\end{solution}
